\chapter{الإحصاء والاحتمال}\label{ch:g9:statproba}

المعطيات والصدفة في كل مكان: نتائج الرياضة، وتوقّعات الطقس،
والألعاب. ويعلّم هذا الفصل كيف تُلخَّص سلسلة أعداد
\omterm{def:g9:statproba:indicators}{بمتوسّطها} \omterm{def:g9:statproba:indicators}{ووسيطها} \omterm{def:g9:statproba:indicators}{ومداها}، وكيف يُحسب \omterm{def:g9:statproba:probability}{اِحتمال} تجارب
عشوائية بسيطة --- بما فيها التجارب ذات الخطوتين، وتُعالَج بأشجار
\omterm{def:g9:statproba:probability}{الاِحتمالات}. ويتوسّع الموضوعان كثيرًا في مجلّد الثانوية من هذه
السلسلة.

\section{تلخيص المعطيات}

\begin{definition}[المتوسّط والوسيط والمدى]\label{def:g9:statproba:indicators}
من أجل سلسلة من $N$ عددًا:
\begin{itemize}
  \item \emph{المتوسّط}\index{متوسط} هو \omterm{def:g1:addition:def}{مجموع} كل القيم مقسومًا
        على $N$؛
  \item \emph{والوسيط}\index{وسيط} قيمة تقسم السلسلة المرتّبة إلى
        نصفين متساويي \omterm{def:g6:measure:volume}{الحجم}: فإذا كان $N$ فرديًا، فهو القيمة
        الوسطى؛ وإذا كان $N$ زوجيًا، فهو منتصف القيمتين
        الوسطيين؛
  \item \emph{والمدى}\index{مدى} هو أكبر قيمة ناقصًا أصغرها.
\end{itemize}
\end{definition}

\begin{example}\label{ex:g9:statproba:indicators}
علامات تلميذ: $8,\ 12,\ 9,\ 15,\ 11$.
\begin{enumerate}
  \item \omterm{def:g9:statproba:indicators}{المتوسّط}: $\dfrac{8 + 12 + 9 + 15 + 11}{5} = \dfrac{55}{5} = 11$.
  \item \omterm{def:g9:statproba:indicators}{والوسيط}: رتّب السلسلة: $8, 9, 11, 12, 15$؛ فالقيمة
        الوسطى (الثالثة) هي $11$.
  \item \omterm{def:g9:statproba:indicators}{والمدى}: $15 - 8 = 7$.
\end{enumerate}
وبعلامة سادسة قدرها $17$: تصير السلسلة المرتّبة
$8, 9, 11, 12, 15, 17$، ويصير \omterm{def:g9:statproba:indicators}{الوسيط} منتصف $11$ و $12$، أي
$11.5$، ويصير \omterm{def:g9:statproba:indicators}{المتوسّط} $\frac{72}{6} = 12$.
\end{example}

\begin{example}[المتوسّط المرجَّح]\label{ex:g9:statproba:weighted}
مقاسات الأحذية المبيعة في يوم:
\begin{center}
\renewcommand{\arraystretch}{1.2}
\begin{tabular}{c|ccccc}
المقاس & $38$ & $39$ & $40$ & $41$ & $42$ \\
\hline
العدد & $3$ & $7$ & $6$ & $3$ & $1$ \\
\end{tabular}
\end{center}
\omterm{def:g9:statproba:indicators}{ومتوسّط} المقاسات هو
\[
\frac{3 \times 38 + 7 \times 39 + 6 \times 40 + 3 \times 41 + 1 \times 42}
{3 + 7 + 6 + 3 + 1}
= \frac{114 + 273 + 240 + 123 + 42}{20}
= \frac{792}{20} = 39.6 .
\]
\end{example}

\begin{remark}
قد يختلف \omterm{def:g9:statproba:indicators}{المتوسّط} \omterm{def:g9:statproba:indicators}{والوسيط} اِختلافًا كبيرًا. ففي السلسلة
$10, 10, 10, 10, 60$ يكون \omterm{def:g9:statproba:indicators}{المتوسّط} $20$ (تجرّه القيمة $60$ إلى
أعلى) بينما \omterm{def:g9:statproba:indicators}{الوسيط} $10$. فاِسأل دائمًا أيّ مؤشّر يمثّل الوضع
تمثيلًا أفضل.
\end{remark}

\section{الاحتمال}

\begin{definition}[الاحتمال]\label{def:g9:statproba:probability}
للتجربة العشوائية عدّة \emph{مخارج} ممكنة؛ و\emph{الحدث} مجموعة
من المخارج. ولكل مخرج \emph{اِحتمال}\index{احتمال}: وهو عدد بين
$0$ و $1$، \omterm{def:g1:addition:def}{ومجموع} اِحتمالات المخارج كلّها $1$؛ واِحتمال $\P(A)$
حدثٍ $A$ هو \omterm{def:g1:addition:def}{مجموع} اِحتمالات مخارجه. وحين تكون المخارج $n$
كلّها متساوية الإمكان،
\[
\P(A) = \frac{\text{عدد المخارج الملائمة}}{n}.
\]
\end{definition}

\begin{example}\label{ex:g9:statproba:spinner}
عجلة مقسومة إلى $8$ قطاعات متساوية: $3$ حمراء، و $3$ زرقاء،
و $2$ خضراء. ولكل قطاع \omterm{def:g9:statproba:probability}{اِحتمال} $\frac18$، إذن
\[
\P(\text{أحمر}) = \frac38, \qquad
\P(\text{أخضر}) = \frac28 = \frac14, \qquad
\P(\text{ليس أخضر}) = 1 - \frac14 = \frac34 .
\]
\end{example}

\begin{proposition}[الحدث المضادّ]\label{prop:g9:statproba:complement}
من أجل كل حدث $A$، \emph{الحدث المضادّ} $\bar A$ («لا يقع $A$»)
يحقّق
\[
\P(\bar A) = 1 - \P(A).
\]
\end{proposition}

\begin{proof}
كل مخرج في واحد بالضبط من $A$ و $\bar A$، إذن
$\P(A) + \P(\bar A)$ يجمع \omterm{def:g9:statproba:probability}{اِحتمالات} كل المخارج، وهو
$1$.
\end{proof}

\section{التجارب ذات الخطوتين}

\begin{method}[أشجار الاحتمالات]\label{met:g9:statproba:tree}
من أجل تجربة في خطوتين:
\begin{enumerate}
  \item اُرسم فرعًا لكل مخرج من الخطوة الأولى، ثم واصِل كل فرع
        بمخارج الخطوة الثانية، كاتبًا كل \omterm{def:g9:statproba:probability}{اِحتمال} على فرعه؛
  \item واِضرب \omterm{def:g9:statproba:probability}{الاِحتمالات} على طول مسار للحصول على \omterm{def:g9:statproba:probability}{اِحتمال} ذلك
        المسار؛
  \item واِجمع \omterm{def:g9:statproba:probability}{اِحتمالات} كل المسارات التي تشكّل حدثًا.
\end{enumerate}
\end{method}

\begin{example}\label{ex:g9:statproba:tree}
كيس فيه $2$ من \omterm{def:g6:lines:objects}{القطع} الحمراء و $3$ سوداء. اِسحب \omterm{def:g6:lines:objects}{قطعة}، وسجّل
لونها، \emph{وأعِدها}، ثم اِسحب من جديد. فيعطي كل سحب أحمر
\omterm{def:g9:statproba:probability}{باِحتمال} $\frac25$، وأسود \omterm{def:g9:statproba:probability}{باِحتمال} $\frac35$.

\begin{omfigure}
\begin{tikzpicture}[scale=0.95]
  \coordinate (root) at (0,0);
  \coordinate (R) at (2.2,1.0);
  \coordinate (B) at (2.2,-1.0);
  \coordinate (RR) at (4.7,1.55);
  \coordinate (RB) at (4.7,0.45);
  \coordinate (BR) at (4.7,-0.45);
  \coordinate (BB) at (4.7,-1.55);
  \draw (root) -- (R) node[midway, above, font=\small, sloped] {$\frac25$};
  \draw (root) -- (B) node[midway, below, font=\small, sloped] {$\frac35$};
  \draw (R) -- (RR) node[midway, above, font=\small, sloped] {$\frac25$};
  \draw (R) -- (RB) node[midway, below, font=\small, sloped] {$\frac35$};
  \draw (B) -- (BR) node[midway, above, font=\small, sloped] {$\frac25$};
  \draw (B) -- (BB) node[midway, below, font=\small, sloped] {$\frac35$};
  \node[omThm, font=\small, above left=-2pt] at (R) {R};
  \node[font=\small, below left=-2pt] at (B) {B};
  \node[right, font=\small] at (RR)
    {RR:\ $\frac25 \times \frac25 = \frac{4}{25}$};
  \node[right, font=\small] at (RB)
    {RB:\ $\frac25 \times \frac35 = \frac{6}{25}$};
  \node[right, font=\small] at (BR)
    {BR:\ $\frac35 \times \frac25 = \frac{6}{25}$};
  \node[right, font=\small] at (BB)
    {BB:\ $\frac35 \times \frac35 = \frac{9}{25}$};
\end{tikzpicture}

{\small شجرة السحبين مع الإعادة: اِضرب على طول الفروع، وتحقّق
أنّ \omterm{def:g1:addition:def}{مجموع} المسارات الأربعة $1$.}
\end{omfigure}

\omterm{def:g9:statproba:probability}{اِحتمال} \omterm{def:g6:lines:objects}{قطعتين} باللون نفسه:
$\frac{4}{25} + \frac{9}{25} = \frac{13}{25}$. \omterm{def:g9:statproba:probability}{واِحتمال} حمراء
واحدة على الأقل: $1 - \P(\text{BB}) = 1 - \frac{9}{25} = \frac{16}{25}$.
\end{example}

\begin{remark}[التواترات تقترب من الاحتمالات]
رمي نرد سليم $6000$ مرّة يعطي \emph{نحو} $1000$ ستّة --- لا
بالضبط. وكلّما زاد عدد التكرارات، اِقتربت التواترات الملاحَظة
أكثر فأكثر من \omterm{def:g9:statproba:probability}{الاِحتمالات}: وهذا ما يجعل \omterm{def:g9:statproba:probability}{الاِحتمال} الأداة
الصحيحة لنمذجة التجارب المكرَّرة (وهي حكاية تُطوَّر في مجلّد
الثانوية من هذه السلسلة).
\end{remark}

\section{تمارين}

\begin{exercise}[$\star$]\label{exo:g9:statproba:1}
اِحسب \omterm{def:g9:statproba:indicators}{المتوسّط} \omterm{def:g9:statproba:indicators}{والوسيط} \omterm{def:g9:statproba:indicators}{والمدى} للسلسلة
$14,\ 9,\ 17,\ 9,\ 12,\ 11,\ 12$.
\end{exercise}

\begin{exercise}[$\star$]\label{exo:g9:statproba:2}
كانت درجات الحرارة عند الظهيرة في أسبوع $19,\ 21,\ 24,\ 18,\ 22,\ 25,\
19$ (بالدرجات). اِحسب \omterm{def:g9:statproba:indicators}{المتوسّط} (إلى العُشر) \omterm{def:g9:statproba:indicators}{والوسيط}.
\end{exercise}

\begin{exercise}[$\star$]\label{exo:g9:statproba:3}
عدد الأهداف التي سجّلها فريق في $10$ مباريات:
\begin{center}
\renewcommand{\arraystretch}{1.2}
\begin{tabular}{c|cccc}
الأهداف & $0$ & $1$ & $2$ & $3$ \\
\hline
المباريات & $2$ & $4$ & $3$ & $1$ \\
\end{tabular}
\end{center}
اِحسب \omterm{def:g9:statproba:indicators}{متوسّط} عدد الأهداف في المباراة، \omterm{def:g9:statproba:indicators}{والوسيط}.
\end{exercise}

\begin{exercise}[$\star$]\label{exo:g9:statproba:4}
يُرمى نرد سليم مرّة واحدة. اِحسب \omterm{def:g9:statproba:probability}{اِحتمال}: «الحصول على
$3$»؛ و«الحصول على \omterm{def:g3:numbers:evenodd}{عدد زوجي}»؛ و«الحصول على $3$ على الأقل»؛
و«عدم الحصول على $6$».
\end{exercise}

\begin{exercise}[$\star$]\label{exo:g9:statproba:5}
علبة فيها $12$ كرة: $5$ بيضاء، و $4$ حمراء، و $3$ خضراء؛ وتُسحب
كرة عشوائيًّا. اِحسب $\P(\text{أبيض})$ و $\P(\text{أحمر أو أخضر})$
و $\P(\text{ليس أبيض})$. فماذا تلاحظ في الأخيرين؟
\end{exercise}

\begin{exercise}[$\star\star$]\label{exo:g9:statproba:6}
تُرمى \omterm{def:g6:lines:objects}{قطعة} نقود، ثم يُرمى نرد سليم. اُرسم الشجرة (أو عُدّ
المخارج) واِحسب \omterm{def:g9:statproba:probability}{اِحتمال} الحصول على صورة \emph{و} $6$؛ ثم \omterm{def:g9:statproba:probability}{اِحتمال}
الحصول على صورة \emph{أو} $6$ (أو كليهما).
\end{exercise}

\begin{exercise}[$\star\star$]\label{exo:g9:statproba:7}
في كيس \cref{ex:g9:statproba:tree} ($2$ حمراء، و $3$ سوداء)،
يُجرى السحبان الآن \emph{من غير} إعادة. اُرسم الشجرة الجديدة
(واِنتبه: \omterm{def:g9:statproba:probability}{اِحتمالات} المستوى الثاني تتغيّر) واِحسب
\omterm{def:g9:statproba:probability}{اِحتمال} سحب \omterm{def:g6:lines:objects}{قطعتين} سوداوين، ثم \omterm{def:g6:lines:objects}{قطعتين} مختلفتَي اللون.
\end{exercise}

\begin{exercise}[$\star\star$]\label{exo:g9:statproba:8}
بعد $9$ مباريات، بلغ \omterm{def:g9:statproba:indicators}{متوسّط} لاعبة كرة سلّة $14$ نقطة في المباراة.
\begin{enumerate}
  \item كم نقطة سجّلت في المجموع؟
  \item وكم نقطة يجب أن تسجّل في المباراة $10$ ليصير \omterm{def:g9:statproba:indicators}{المتوسّط}
        $15$؟
\end{enumerate}
\end{exercise}

\begin{exercise}[$\star\star\star$]\label{exo:g9:statproba:9}
لعبة: اِرمِ نردين سليمين؛ وتفوز إذا تساوت النتيجتان
(«الزوج»).
\begin{enumerate}
  \item اِحسب \omterm{def:g9:statproba:probability}{اِحتمال} الفوز.
  \item وتلعب مرّتين (لعبتان مستقلّتان). وباِستعمال شجرة، اِحسب
        \omterm{def:g9:statproba:probability}{اِحتمال} الفوز مرّة واحدة على الأقل.
\end{enumerate}
\end{exercise}

\section{مسألة: رهانات الفارس المهلِكة}

\begin{problem}[{مسألة نهاية الأسبوع --- نزاع القمار سنة 1654 الذي
أنشأ نظرية الاحتمال، ومصادفة أعياد الميلاد التي تخدع كل قسم}]
\label{pb:g9:statproba:1}
في سنة 1654، أثرى نبيل فرنسي مولع بالقمار، هو الفارس دي ميري،
من رهان نرد واحد، ثم بدأ يخسر في رهان آخر ظنّه مكافئًا له.
وحائرًا، سأل عالم الرياضيات بليز باسكال؛ فكتب باسكال إلى بيير دو
فيرما؛ وأسّس تبادل رسائلهما نظرية \omterm{def:g9:statproba:probability}{الاحتمال}. وتعيد هذه المسألة
تمثيل القضية كلّها بأدوات هذا الفصل --- المخارج المتساوية
الإمكان، والأشجار، وقاعدة الحدث المضادّ
(\cref{prop:g9:statproba:complement}) --- وتنتهي بأكثر المصادفات
مخالفةً للحدس على الإطلاق.

\medskip
\noindent\textbf{الجزء الأول --- النرد، معدودًا بأمانة.}
\begin{enumerate}
  \item اِرمِ نردًا سليمًا مرّتين. وباِستعمال شجرة (أو جدول
        $6 \times 6$)، اِحسب \omterm{def:g9:statproba:probability}{اِحتمال} \emph{عدم} ظهور ستّة في
        الرميتين، واِستنتج \omterm{def:g9:statproba:probability}{اِحتمال} ظهور ستّة واحدة على الأقل.
  \item ويجادل زميل: «فرصة واحدة من ستّ في كل رمية، ومرّتين،
        إذن $\frac16 + \frac16 = \frac13$». قارن بالسؤال 1
        وأشِر إلى المخارج التي يعدّها جمعُه مرّتين بالضبط.
  \item وبرمي نردين وجمع النتيجتين: اِحسب \omterm{def:g9:statproba:probability}{اِحتمال} \omterm{def:g1:addition:def}{مجموع}
        قدره $7$ \omterm{def:g1:addition:def}{ومجموع} قدره $2$. ولماذا يكون $7$ مفضَّل
        المقامرين؟
  \item ونزاع قديم: بنردين، أيّهما أرجح \omterm{def:g1:addition:def}{مجموع} قدره $9$ أم
        \omterm{def:g1:addition:def}{مجموع} قدره $10$؟ عُدّ المخارج واحسم الأمر.
  \item وعمّم السؤال 1: فسّر لماذا يكون \omterm{def:g9:statproba:probability}{اِحتمال} ظهور ستّة
        واحدة على الأقل في $n$ رمية هو
        \[
        1 - \left(\frac56\right)^{\!n} ,
        \]
        إذ تتضاعف فروع «لا ستّة» في الشجرة من مستوى إلى مستوى.
\end{enumerate}

\medskip
\noindent\textbf{الجزء الثاني --- رهانا الفارس.}
\begin{enumerate}[resume]
  \item الرهان الأول، بمال متكافئ: \emph{ستّة واحدة على
        الأقل في أربع رميات نرد}. اِحسب
        $\left(\frac56\right)^4$ كسرًا وعددًا عشريًّا،
        \omterm{def:g9:statproba:probability}{واِحتمال} فوز الفارس. فهل كان الرهان جيدًا؟
  \item وكان تعليل الفارس نفسه: «أربع رميات، كل واحدة
        $\frac16$: أي $\frac46$ من الفرصة». وقد أعطى جوابًا
        يكاد يكون صحيحًا هنا --- لكن اِدفعه إلى $7$ رميات: فأيّ
        عبث ينتجه، وأيّ خطأ من السؤال 2 يكرّره؟
  \item والرهان الثاني، الذي ظنّه مكافئًا \omterm{def:g9:linfunc:linear}{بالتناسب}
        («$24$ رمية، كل واحدة $\frac1{36}$: أي من جديد
        $\frac{24}{36} = \frac46$»): \emph{ستّة مزدوجة واحدة
        على الأقل في $24$ رمية نردين}. اِحسب \omterm{def:g9:statproba:probability}{اِحتمال} الفوز
        الحقيقي $1 - \left(\frac{35}{36}\right)^{24}$
        (بالآلة الحاسبة). ولماذا خرب الفارس ببطء؟
  \item وكم رمية نردين كانت ستعيد إليه أفضليته؟ اِختبر $n = 25$
        بالآلة الحاسبة واِستنتج.
  \item واُذكر العبرة في جملتين: ما الخطأ، في العموم، في ضرب
        \omterm{def:g9:statproba:probability}{اِحتمال} في عدد المحاولات --- وأيّ أداة صحيحة تحلّ محلّ
        هذا الإغراء؟
\end{enumerate}

\medskip
\noindent\textbf{الجزء الثالث --- \omterm{def:g9:statproba:indicators}{المتوسّطات} \omterm{def:g9:statproba:indicators}{والوسطاء} وأعياد الميلاد.}
\begin{enumerate}[resume]
  \item شركة صغيرة تدفع لتسعة موظّفين $2\,000$ يورو في الشهر
        وللمدير $20\,000$. اِحسب \omterm{def:g9:statproba:indicators}{المتوسّط} \omterm{def:g9:statproba:indicators}{والوسيط} للأجور
        (\cref{def:g9:statproba:indicators}). فأيّ عدد ينبغي أن
        يذكره إعلان التوظيف بأمانة؟
  \item وأنشِئ قائمة من خمس علامات \omterm{def:g9:statproba:indicators}{متوسّطها} $12$
        \omterm{def:g9:statproba:indicators}{ووسيطها} $15$. فماذا يقول هذا الزوج (\omterm{def:g9:statproba:indicators}{متوسّط} دون
        \omterm{def:g9:statproba:indicators}{الوسيط}) عن شكل العلامات؟
  \item ومسألة أعياد الميلاد: في قسم فيه $23$ تلميذًا، ما
        \omterm{def:g9:statproba:probability}{اِحتمال} أن يتشارك اِثنان على الأقل عيد ميلاد؟ كوّن
        الحدث المضادّ (أعياد الميلاد $23$ كلّها مختلفة):
        \[
        \frac{365}{365} \times \frac{364}{365} \times
        \frac{363}{365} \times \dots \times \frac{343}{365} .
        \]
        اِحسب \omterm{def:g2:mult:def}{جداء} العوامل الثلاثة الأولى، وصِف كيف تتطوّر
        العوامل، ثم --- بما أنّ \omterm{def:g2:mult:def}{الجداء} الكامل نحو $0.493$ ---
        أجِب عن السؤال. ويتوقّع معظم الناس «مستبعَد جدًّا»:
        فماذا تقول الرياضيات؟
  \item والحدس مصلَّحًا: كم \emph{زوجًا} من التلاميذ يمكن أن
        يتشارك عيد ميلاد في قسم فيه $23$ تلميذًا (عدّ
        المصافحات في \cref{pb:g6:wholes:1})؟ فسّر في جملة واحدة
        لماذا يقود هذا العدد المفاجأة، لا العدد $23$ نفسه.
  \item والختام: صِف تجربة للتحقّق من إحدى النتيجتين --- رهان
        دي ميري بنرد حقيقي، أو مسألة أعياد الميلاد عبر عدّة
        أقسام --- واُذكر بعناية ما تتوقّعه من التواترات
        الملاحَظة كلّما زاد عدد التكرارات. (ولذلك التوقّع
        اِسم، هو قانون الأعداد الكبيرة؛ وبرهانه ينتظر في مجلّد
        الثانوية.)
\end{enumerate}
\end{problem}
